% !PRINTED 290
\weberpage{290}

Setzt man zwei dieser Substitutionen zusammen, so hat man zu beachten, dass
\[
(\psi, \psi_b) = (\psi_a, \psi_{ab})
\]
ist, so dass man
\[
(\psi, \psi_a)(\psi, \psi_b) = (\psi, \psi_{ab})
\]
hat. Es setzen sich also diese Substitutionen ganz in derselben Weise zusammen, wie
nach \S.\,4 dieses Bandes die Nebengruppen, und es ergiebt sich daraus:

\noindent 1. \gesperrt{Die Galois'sche Gruppe der zu $Q$ geh\"origen Partialresolvente
ist isomorph mit der zu $Q$ complement\"aren Gruppe $P/Q$.}

Nehmen wir jetzt irgend eine Compositionsreihe von $P$ mit der zugeh\"origen
Indexreihe (\S.\,6):
\weq{1}{
\begin{gathered}
P,\ P_1,\ P_2,\ \ldots,\ P_{\mu-1},\ 1 \\
j_1,\ j_2,\ \ldots,\ j_{\mu-1},\ j_\mu,
\end{gathered}
}
so wird nach dem soeben Gesagten die Gruppe der Gleichung $f = 0$ von $P$ auf $P_1$
reducirt durch Adjunction \gesperrt{einer} Wurzel einer Normalgleichung vom Grade $j_1$. Dann
wird sie auf $P_2$ reducirt durch eine Wurzel einer Normalgleichung vom Grade $j_2$
u.\,s.\,f. und endlich wird die Gleichung vollst\"andig gel\"ost durch eine Wurzel
einer Normalgleichung vom Grade $j_\mu$.

Auf dies Aufl\"osungsverfahren f\"allt nun von dem Satze \"uber die
Unver\"anderlichkeit der Indexreihe (\S.\,6, I.) ein neues Licht.

\noindent 2. \gesperrt{Um die Gleichung $f = 0$ zu l\"osen, hat man nach einander je
eine Wurzel einer Normalgleichung der Grade $j_1, j_2, \ldots, j_\mu$ zu adjungiren.
Die Grade dieser Resolventen k\"onnen zwar in der Reihenfolge, nicht aber in der
Gesammtheit abge\"andert werden.}

Die Zahlen $j_1, j_2, \ldots, j_\mu$ h\"angen also weit tiefer mit der Natur einer
Gleichung oder allgemeiner mit den durch die Gleichung definirten K\"orpern zusammen,
als etwa der Grad der Gleichung; denn w\"ahrend der Grad durch Transformation auf
mannigfache Weise ver\"andert werden kann, wenn man Functionen der Wurzeln als neue
Unbekannte einf\"uhrt, bleiben die Zahlen $j_1, j_2, \ldots, j_\mu$ immer erhalten und
sind als wahre \gesperrt{Invarianten} des durch die Gleichung definirten
Normalk\"orpers zu betrachten.
